\chapter{Results}
\label{chap:results}

This chapter reports monthly local-projection responses to a one-standard-deviation ECB target-factor easing surprise. The evidentiary object is the response path: sign, peak horizon, accumulated movement, stability, and uncertainty. The chapter follows the interpretation hierarchy of the thesis: identification validity first, housing-credit transmission second, intermediation variables third, financial-market responses fourth, compensation-linked variables fifth, and regime comparisons last.

\section{Identification Validity}

The baseline shock is \nolinkurl{target_factor_monthly_easing}, a monthly aggregation of high-frequency ECB target-factor surprises signed so positive values denote easing. The strongest low-frequency bridge is the end-of-month DFR change, with \(F=17.07\) and partial \(R^2=0.066\). This provides meaningful full-sample relevance, but the jackknife and rolling diagnostics show sensitivity to event composition. The interpretation therefore remains policy-news based rather than a stronger structural treatment claim.

Information effects are a second identification boundary. ECB announcements can convey policy accommodation and central-bank information about future macroeconomic conditions at the same time. The clean-event and contaminated-event diagnostics are therefore used to check whether the strongest response patterns survive possible announcement contamination. This boundary is especially important for equities, where the baseline real DAX response is negative after an easing surprise.

\artifacttable
  {Instrument strength matrix}
  {tab:results-instrument-strength-matrix}
  {../results/final/diagnostics/instrument_strength_matrix.csv}
  {Diagnostic matrix used to evaluate relevance and stability of candidate monetary surprise measures.}

\artifacttable
  {Information-effect event screen}
  {tab:results-information-effect-screen}
  {../results/final/diagnostics/information_effect_event_screen.csv}
  {Event-level screen used to separate clean events from information-effect-sensitive observations.}

\thesisfigure[0.98\textwidth]
  {enhancement/comparative_irf_panel.png}
  {Comparative impulse-response panels for housing-credit, lending-condition, market-liquidity, and compensation-labor variables.}
  {fig:comparative-irf-panel}

\Cref{fig:comparative-irf-panel} previews the main empirical pattern. House-purchase lending growth is positive through h12 and accumulates durably. Mortgage and NFC spreads remain positive over the response path. The wage tracker is weaker and loses long-horizon precision, while employment expectations reverse by h24. The DAX moves in the opposite direction, reinforcing the need for information-effect interpretation.

\section{Housing-Credit Transmission}

Housing-credit transmission is the core result. The clearest and most persistent response is house-purchase lending growth.

\thesisfigure
  {baseline/monthly_irf_ecb_house_purchase_growth_yoy.png}
  {Baseline impulse response of house-purchase lending growth to a one-standard-deviation ECB easing surprise.}
  {fig:baseline-house-purchase-growth}

\thesisfigure
  {cumulative/cumulative_interval_ecb_house_purchase_growth_yoy.png}
  {Cumulative response interval for house-purchase lending growth.}
  {fig:cumulative-house-purchase-growth}

House-purchase lending growth is positive from h0 through h12, peaks at h6, and remains positive in accumulated terms through h24. This makes it the centerpiece variable of the thesis. The response is not interpreted as a direct house-price effect or as household welfare evidence. It is reduced-form evidence that ECB easing news leaves a strong and persistent trace in the housing-credit channel.

The stock-flow distinction is essential. House-purchase lending growth is a stock-based credit-growth measure. It captures cumulative financing conditions, sustained transmission, and medium-run credit accommodation. Pure-new house-purchase loans are a flow-based origination measure. They capture immediate lending behavior and are more volatile, shorter in sample, and more sensitive to month-specific origination dynamics.

\thesisfigure
  {baseline/monthly_irf_ln_ecb_house_purchase_pure_new_loans.png}
  {Baseline impulse response of pure new house-purchase loans to a one-standard-deviation ECB easing surprise.}
  {fig:baseline-house-purchase-loans}

\thesisfigure
  {cumulative/cumulative_interval_ln_ecb_house_purchase_pure_new_loans.png}
  {Cumulative response interval for pure new house-purchase loans.}
  {fig:cumulative-house-purchase-loans}

Pure-new house-purchase loans support housing-credit sensitivity at short horizons but do not replace the canonical annual-growth proxy. The divergence between the growth stock and new-flow measures is therefore not a failure of the mechanism. It shows that sustained credit accommodation is more visible in stock-based house-purchase lending growth than in immediate monthly originations.

\section{Lending Conditions and Intermediation}

The intermediation layer asks whether lending conditions and credit variables make the housing-credit response economically plausible. Mortgage and NFC lending spreads provide the cleanest evidence.

\thesisfigure
  {baseline/monthly_irf_ecb_mir_mortgage_lending_spread_dfr.png}
  {Baseline impulse response of the mortgage lending spread to a one-standard-deviation ECB easing surprise.}
  {fig:baseline-mortgage-spread}

\thesisfigure
  {baseline/monthly_irf_ecb_mir_nfc_lending_spread_dfr.png}
  {Baseline impulse response of the NFC lending spread to a one-standard-deviation ECB easing surprise.}
  {fig:baseline-nfc-spread}

\thesisfigure
  {cumulative/cumulative_interval_ecb_mir_mortgage_lending_spread_dfr.png}
  {Cumulative response interval for the mortgage lending spread.}
  {fig:cumulative-mortgage-spread}

The mortgage spread is positive at every reported horizon, peaks at h12, and has positive h6, h12, and h24 accumulated responses whose 90 percent intervals exclude zero. The NFC spread follows the same profile at smaller magnitude. Lending-rate levels are less precise, so the result rests on margins relative to the DFR rather than on a mechanical rate-level story.

Broad credit quantities are weaker and less central than the housing-credit proxy.

\thesisfigure
  {baseline/monthly_irf_ln_nfc_loans_ea_stock.png}
  {Baseline impulse response of Euro Area NFC loan stocks to a one-standard-deviation ECB easing surprise.}
  {fig:baseline-nfc-loans}

NFC and household credit stocks have smaller paths and accumulated intervals that include zero. This contrast is important: the thesis does not claim that all credit aggregates respond strongly. It claims that transmission is sectorally differentiated, with housing-credit growth and lending margins showing stronger relative persistence than broad productive-credit quantities.

The sequential timing outputs support the intermediation reading without identifying structural mediation.

\thesisfigure
  {mechanism/sequential_timing_response_map.png}
  {Sequential timing map for banking, credit, housing-credit, and wage-pressure pathways.}
  {fig:results-sequential-timing-map}

\artifacttable
  {Sequential transmission outputs}
  {tab:sequential-transmission}
  {../results/final/mechanism/sequential_timing_outputs.csv}
  {Ordered timing associations linking lending conditions, credit quantities, housing credit, and wage-pressure proxies.}

\textbf{Observation.} The sequential outputs order mortgage spreads, household credit, house-purchase growth, new mortgage flows, NFC spreads, NFC credit, and liquidity variables. \textbf{Interpretive limit.} The sequencing aligns with intermediary propagation. \textbf{Identification boundary.} It remains ordered timing evidence because the variables respond to the same policy-news shock.

\section{Equity and Asset-Market Responses}

Financial-market responses are mixed. They do not support a simple claim that easing mechanically inflates all asset prices.

\thesisfigure
  {baseline/monthly_irf_ln_dax_real_de.png}
  {Baseline impulse response of the real DAX to a one-standard-deviation ECB easing surprise.}
  {fig:baseline-dax}

\thesisfigure
  {cumulative/cumulative_interval_ln_dax_real_de.png}
  {Cumulative response interval for the real DAX.}
  {fig:cumulative-dax}

The real DAX is small at h0--h1 and increasingly negative from h3 to h24; h12 and h24 intervals exclude zero. This is the largest sign-coherence issue in the financial block, and it is addressed directly. A negative equity response after an easing surprise may indicate that the announcement contains adverse macroeconomic information as well as accommodative policy news. Markets react to both policy and information. Easing can therefore coincide with deteriorating growth expectations, weaker expected earnings, or higher risk premia.

\thesisfigure
  {baseline/monthly_irf_ln_ecb_assets_ea_stock.png}
  {Baseline impulse response of ECB assets to a one-standard-deviation ECB easing surprise.}
  {fig:baseline-ecb-assets}

ECB assets are also negative through h24 but imprecise. They provide liquidity-stock context rather than a strong response pillar. The asset-market evidence is therefore a boundary condition: selected financial variables respond, but the housing-credit and lending-condition blocks carry the stronger mechanism claim.

\section{Compensation-Linked and Real-Economy Responses}

The compensation-linked block provides the comparative real-economy evidence. It is not silent, but it is weaker and less consistent than the housing-credit and intermediation evidence.

\thesisfigure
  {baseline/monthly_irf_ecb_wage_tracker_ex_oneoffs_real_yoy.png}
  {Baseline impulse response of the real ECB Wage Tracker excluding one-offs to a one-standard-deviation ECB easing surprise.}
  {fig:baseline-wage-tracker}

\thesisfigure
  {cumulative/cumulative_interval_ecb_wage_tracker_ex_oneoffs_real_yoy.png}
  {Cumulative response interval for the real wage tracker excluding one-offs.}
  {fig:cumulative-wage-tracker}

\thesisfigure
  {baseline/monthly_irf_eurostat_ecfin_eei_ea20.png}
  {Baseline impulse response of employment expectations to a one-standard-deviation ECB easing surprise.}
  {fig:baseline-employment-expectations}

The real ECB Wage Tracker excluding one-offs is negative from h0 through h12 and imprecisely positive at h24. Its accumulated response excludes zero at h6 and h12 but not h24. Employment expectations move earlier and reverse at h24. The German industry wage-bill proxy is more positive at medium horizons, but remains robustness evidence because it is sectoral and national. Retail volume is mostly weak. The comparative conclusion is therefore that broad real-income transmission is less consistent than housing-credit transmission, not that labor-market channels are absent.

\section{Persistence Synthesis}

The strongest claims occur where sign, accumulated response, uncertainty, and stability align.

\begin{table}[htbp]
  \centering
  \caption{Master synthesis table for persistence evidence}
  \label{tab:persistence-synthesis}
  \scriptsize
  \setlength{\tabcolsep}{3pt}
  \begin{tabularx}{\textwidth}{@{}>{\raggedright\arraybackslash}p{0.16\textwidth}>{\raggedright\arraybackslash}p{0.10\textwidth}>{\raggedright\arraybackslash}p{0.08\textwidth}>{\raggedright\arraybackslash}p{0.16\textwidth}>{\raggedright\arraybackslash}p{0.16\textwidth}>{\raggedright\arraybackslash}X@{}}
    \toprule
    Variable & Sign & Peak horizon & Cumulative persistence & Stability & Interpretation \\
    \midrule
    House-purchase lending growth & Positive & h6 & Positive h24; 90\% CI excludes zero & High: sign 0.86; cumulative 1.00 & Core housing-credit response and main empirical contribution. \\
    Mortgage spreads & Positive & h12 & Positive h24; 90\% CI excludes zero & Very high: sign 1.00; cumulative 1.00 & Stable lending-condition margin supporting intermediation. \\
    NFC spreads & Positive & h12 & Positive h24; 90\% CI excludes zero & High: sign 0.90; cumulative 1.00 & Smaller but durable corporate credit-price adjustment. \\
    NFC credit & Mixed & h24 & Negative h24; CI crosses zero & Weak: sign 0.52; cumulative 0.50 & Productive-credit quantity response too imprecise for a persistence claim. \\
    DAX & Negative & h24 & Negative h24; 90\% CI excludes zero & High: sign 0.90; cumulative 1.00 & Information-sensitive market response, not asset-appreciation evidence. \\
    ECB assets & Negative & h24 & Negative h24; CI crosses zero & Moderate: sign 0.79; cumulative 0.86 & Liquidity-stock context, not a strong persistence pillar. \\
    Wage tracker & Mixed & h24 & Negative h24; CI crosses zero & Moderate: sign 0.67; cumulative 0.86 & Compensation-linked response attenuates by h24. \\
    Employment expectations & Positive then reverse & h1 & Positive h24; CI crosses zero & Moderate: sign 0.76; cumulative 0.79 & Short-horizon labor expectations, not durable income transmission. \\
    \bottomrule
  \end{tabularx}
  \vspace{0.25em}
  {\footnotesize Source: \nolinkurl{results/final/tables/peak_response_table.csv}, \nolinkurl{results/final/tables/cumulative_persistence_table.csv}, \nolinkurl{results/final/tables/stability_metrics.csv}, and \nolinkurl{results/final/uncertainty/significance_heatmap.csv}.}
\end{table}

\Cref{tab:persistence-synthesis} is the inferential map. House-purchase lending growth and lending spreads carry the strongest positive horizon durability. The DAX is sustained in the opposite direction and therefore belongs in the information-effects boundary. ECB assets, wage pressure, employment expectations, and NFC credit mark limits because precision, sign stability, or accumulated exclusion weakens.

\thesisfigure[0.95\textwidth]
  {enhancement/sign_persistence_map.png}
  {Sign persistence map across baseline local-projection horizons.}
  {fig:sign-persistence-map}

\thesisfigure[0.95\textwidth]
  {enhancement/horizon_durability_matrix.png}
  {Cumulative horizon durability matrix showing where 90 percent cumulative intervals exclude zero.}
  {fig:horizon-durability-matrix}

\artifacttable
  {Normalized baseline IRF output}
  {tab:normalized-irf-output}
  {../results/final/tables/normalized_irf_outputs.csv}
  {Canonical horizon-level response estimates, intervals, bootstrap bands, and shock scaling used to populate thesis results.}

\artifacttable
  {Horizon significance table}
  {tab:horizon-significance}
  {../results/final/tables/horizon_significance_table.csv}
  {Horizon-by-response significance classifications for the baseline target-factor shock.}

\artifacttable
  {Cumulative persistence table}
  {tab:cumulative-persistence}
  {../results/final/tables/cumulative_persistence_table.csv}
  {Canonical persistence comparison used for financial, housing-credit, banking-spread, and compensation-linked claims.}

\section{Hypothesis Assessment}

\begin{table}[htbp]
  \centering
  \caption{Empirical assessment of the three thesis hypotheses}
  \label{tab:hypothesis-assessment}
  \small
  \setlength{\tabcolsep}{5pt}
  \begin{tabularx}{\textwidth}{@{}>{\raggedright\arraybackslash}p{0.24\textwidth}>{\raggedright\arraybackslash}X>{\raggedright\arraybackslash}p{0.20\textwidth}@{}}
    \toprule
    Hypothesis & Evidence & Assessment \\
    \midrule
    H1: Housing-credit persistence & House-purchase lending growth is positive, peaks at h6, and remains cumulatively durable through h24. Pure-new mortgage flows support short-horizon sensitivity but are more volatile and shorter in sample. & Supported. \\
    H2: Compensation-linked attenuation & The wage tracker moves at short and medium horizons but loses h24 cumulative exclusion, while employment expectations reverse and broader real-activity evidence is weak. & Supported with proxy boundary. \\
    H3: Differential financial-intermediation transmission & Housing-credit growth and lending spreads are stronger and more persistent than NFC credit quantities and compensation-linked proxies. The sequential timing outputs align with intermediation but remain reduced-form timing evidence. & Conditionally supported. \\
    \bottomrule
  \end{tabularx}
\end{table}

\Cref{tab:hypothesis-assessment} addresses the three hypotheses directly. The evidence supports a housing-credit persistence result, supports weaker compensation-linked transmission within proxy limits, and conditionally supports differential financial-intermediation transmission. The conditionality matters: the thesis' strongest claim is not that every coefficient is invariant, but that the ordering of durable responses survives most clearly for house-purchase lending growth and lending spreads.

\section{FEVD Interpretation}

The final pipeline does not export a thesis-facing \(\FEVD\) decomposition under \nolinkurl{results/final/}. The chapter therefore keeps the focus on horizon responses and accumulated paths.

\section{Regime Comparisons}

Regime comparisons are descriptive.

\thesisfigure
  {regimes/regime_h6_response_comparison.png}
  {Six-month response comparison across regimes.}
  {fig:regime-h6-comparison}

Pre-QE estimates show strong positive accumulated house-purchase lending growth through h24. QE maps the target surprise into different observed conditions, including negative house-purchase lending growth and negative NFC spreads. COVID has short-horizon visibility only. Tightening shows positive housing-credit and spread accumulation with negative wage-pressure cumulation. These patterns inform historical dependence; they do not replace the full-sample hierarchy.

\artifacttable
  {Regime-specific reduced-form LP output}
  {tab:regime-lp-output}
  {../results/final/regime/monthly_regime_reduced_form_lp.csv}
  {Descriptive regime-specific response estimates used to assess historical dependence.}
